%% kapa: exam=2 round=28 year=2017 period=1 subject=appraisal_practice subject_ko=감정평가실무 seq=4 number=4 points=10
%% source: Q-Net 감정평가사 기출문제 (cst00309) artlSeq=5201712 file=raw/1008/5201712/2017년 감정평가사 2차시험 기출문제(감정평가 실무).pdf
%% conversion: kordoc 4.24.51 -> tools/gen2.py (auto); review=proofread-sonnet; flags=none
\begin{longquestion}{【문제 4】}
감정평가사 김○○은 다음 물건에 대하여 A기업으로부터 일반거래시가참고용 감정평가를 의뢰받았다. 기준시점을 2017.07.01로 하여 관련 법규 및 이론을 참작하고 제시된 자료를 활용하여 감정평가하시오. (10점) \par
<자료 1> 대상물건 개요 \par
{\small\centering
\begin{tabular}{|c|c|c|c|c|c|c|}
\hline
\makecell{소재지} & \multicolumn{6}{|c|}{\makecell{A시 B구 C동 100외 2필지}} \\
\hline
\makecell{건물명, 층, 호수} & \multicolumn{6}{|c|}{\makecell{“D타운” 제10동 제17층 제1706호}} \\
\hline
\makecell{용도} & \multicolumn{2}{|c|}{\makecell{아파트}} & \multicolumn{2}{|c|}{\makecell{사용승인일}} & \multicolumn{2}{|c|}{\makecell{2010.10.10}} \\
\hline
\multirow{2}{*}{\makecell{면적}} & \makecell{전유면적(m\textsuperscript{2})} & \multicolumn{2}{|c|}{\makecell{공용면적(m\textsuperscript{2})}} & \multicolumn{2}{|c|}{\makecell{분양면적(m\textsuperscript{2})}} & \makecell{대지권면적(m\textsuperscript{2})} \\
\cline{2-7}
 & \makecell{85} & \multicolumn{2}{|c|}{\makecell{25}} & \multicolumn{2}{|c|}{\makecell{110}} & \makecell{28.5} \\
\hline
\end{tabular}\par}

<자료 2> 현장조사내용 \par
1. 본건 인근은 아파트, 다세대 및 다가구주택, 상가, 공장 등이 혼재하는 지역임 \par
2. 본건까지 차량출입 가능하고 인근에 노선버스정류장이 위치해 있음 \par
3. 본건은 현재 방 3개, 주방, 거실, 화장실 2개, 발코니로 구성되어 있으며, 위생 및 급ㆍ배수설비, 난방설비, 승강기설비, 소화전설비 등이 되어 있음 \par
4. 본건은 발코니가 합법적으로 확장되어 있으며 확장면적은 10m\textsuperscript{2}임 \par
5. 본건의 관리상태를 상, 중, 하로 나타낼 경우 ‘하’에 해당함 \par
6. 대상물건이 위치한 동일 아파트단지 내의 거래사례를 분석한 결과, 거래시점, 발코니 확장정도, 관리상태, 층, 향, 동에 따라 가격차이가 존재함 \par
<자료 3> 인근 아파트 거래사례(정상적인 거래로 판단됨) \par
{\small\centering
\begin{tabular}{|c|c|c|c|c|c|c|c|c|}
\hline
\makecell{기호} & \makecell{소재지,\newline{}지번, 명칭} & \makecell{동} & \makecell{층} & \makecell{호수} & \makecell{전유면적\newline{}(m\textsuperscript{2})} & \makecell{거래금액(원)} & \makecell{거래시점} & \makecell{관리상태} \\
\hline
\makecell{1} & \multirow{5}{*}{\makecell{A시 B구\newline{}C동 100외\newline{}2필지\newline{}D타운}} & \makecell{제10동} & \makecell{4} & \makecell{406} & \makecell{85} & \makecell{338,750,000} & \makecell{2016.06.10} & \makecell{하} \\
\cline{1-1} \cline{3-9}
\makecell{2} &  & \makecell{제10동} & \makecell{18} & \makecell{1803} & \makecell{85} & \makecell{335,000,000} & \makecell{2016.12.10} & \makecell{하} \\
\cline{1-1} \cline{3-9}
\makecell{3} &  & \makecell{제10동} & \makecell{16} & \makecell{1605} & \makecell{85} & \makecell{350,000,000} & \makecell{2017.03.25} & \makecell{중} \\
\cline{1-1} \cline{3-9}
\makecell{4} &  & \makecell{제15동} & \makecell{8} & \makecell{802} & \makecell{85} & \makecell{338,000,000} & \makecell{2017.07.01} & \makecell{상} \\
\cline{1-1} \cline{3-9}
\makecell{5} &  & \makecell{제15동} & \makecell{8} & \makecell{801} & \makecell{85} & \makecell{345,000,000} & \makecell{2017.07.01} & \makecell{중} \\
\hline
\end{tabular}\par}

<자료 4> A시 월별 아파트 매매가격지수 \par
{\small\centering
\begin{tabular}{|c|c|c|c|c|c|c|c|c|c|c|c|c|}
\hline
\multicolumn{6}{|c|}{\makecell{2016년}} & \multicolumn{7}{|c|}{\makecell{2017년}} \\
\hline
\makecell{7월} & \makecell{8월} & \makecell{9월} & \makecell{10월} & \makecell{11월} & \makecell{12월} & \makecell{1월} & \makecell{2월} & \makecell{3월} & \makecell{4월} & \makecell{5월} & \makecell{6월} & \makecell{7월} \\
\hline
\makecell{103.4} & \makecell{103.6} & \makecell{103.8} & \makecell{104} & \makecell{104.4} & \makecell{104.4} & \makecell{104.4} & \makecell{104.4} & \makecell{104.5} & \makecell{104.6} & \makecell{104.7} & \makecell{105.0} & \makecell{105.1} \\
\hline
\end{tabular}\par}

<자료 5> 기타자료 \par
1. 거래사례 중 기호5는 발코니가 합법적으로 10m\textsuperscript{2} 확장된 것으로 조사되었고 나머지는 확장되지 않은 것으로 조사됨 \par
2. 본건과 거래사례는 방배치 등 기타 구조 측면에서 모두 동일함 \par
3. 관리상태에 따른 가격격차 정도는 다음과 같음(D타운 전체 적용가능) \par
{\small\centering
\begin{tabular}{|c|c|c|}
\hline
\makecell{하} & \makecell{중} & \makecell{상} \\
\hline
\makecell{100} & \makecell{101} & \makecell{102} \\
\hline
\end{tabular}\par}

4. 층에 따른 가격격차 정도는 다음과 같음(D타운 전체 적용가능) \par
{\small\centering
\begin{tabular}{|c|c|c|}
\hline
\makecell{1층 ∼ 3층} & \makecell{4층 ∼ 10층} & \makecell{11층 ∼ 20층} \\
\hline
\makecell{100} & \makecell{105} & \makecell{108} \\
\hline
\end{tabular}\par}

5. 본건 아파트의 1호 ∼ 3호는 남동향이며, 4호 ∼ 6호는 남향임 \par
\end{longquestion}
